
Multi-dimensional trustworthiness evaluation can miss compound failures invisible to independent assessment. A medical diagnosis LLM might achieve 92\% truthfulness and 89\% fairness when evaluated separately, yet fail on both dimensions for the same patients---a coupling pattern current benchmarks cannot detect.

Through phi coefficient analysis demonstrated on synthetic benchmark emulation across 5 trustworthiness dimensions and 3 simulated model profiles, this study establishes that coupling is detectable when designed into data (phi 0.33-0.40, p < 1e-13) but measurably sparse---limited to 2 dominant dimension pairs, not pervasive across all 10 possible pairs. Truthfulness-robustness coupling (phi 0.36-0.40) and fairness-safety coupling (phi 0.33-0.40) appear in synthetic validation. No model exhibits coupling in 3 or more pairs after correction for multiple comparisons.

This sparsity detection capability persists when controlling for instance difficulty. Partial correlation analysis yields partial phi values of 0.36-0.56 with retention of 86-161\%---in five of six cases, controlling for difficulty strengthens rather than weakens coupling in synthetic data by design. The methodology isolates designed coupling from confounds.

Models show distinct coupling profiles in synthetic data---GPT-4 shows truthfulness-robustness chain, Claude-3 shows fairness-safety-privacy cluster, Llama-3 shows minimal coupling---though statistical confirmation requires larger samples than the proof-of-concept provides.

All findings derive from synthetic coupling data validating measurement methodology. Real-world coupling magnitudes in production LLMs remain uncharacterized pending access to MultiTrust and TrustLLM benchmarks. If validated on real models, the implications are immediate: multi-dimensional evaluation frameworks should report coupling statistics alongside per-dimension scores, and model selection teams could exploit coupling patterns when deployment scenarios demand multiple trustworthiness guarantees simultaneously.

Trustworthiness dimensions are neither universally coupled nor universally independent. The methodology demonstrates they can exhibit sparse coupling, a detectable property requiring targeted mitigation strategies rather than universal solutions. Real-world magnitudes remain unknown pending production benchmark access.
